\documentclass[border=5mm]{standalone}
% ==============================================================================
% SETUP.TEX - CONFIGURACIÓN GLOBAL DEL PROYECTO
% ==============================================================================
% Este archivo centraliza todos los paquetes y estilos.
% Debe incluirse en el preámbulo del libro principal y de los archivos standalone.
% \PassOptionsToPackage{gray}{xcolor}
% --- 1. CODIFICACIÓN E IDIOMA ---
% fix-cm habilita las fuentes Computer Modern en tamaños arbitrarios (por debajo
% de 5 pt, entre otros). Debe cargarse antes que fontenc.
\usepackage{fix-cm}
\usepackage[utf8]{inputenc}
\usepackage[T1]{fontenc}
\usepackage[spanish, es-tabla]{babel}  
\usepackage{csquotes} % Recomendado para babel y citas

\usepackage{import}
% mode=image: Usa el PDF pre-compilado, nunca intenta recompilar.
% Debes compilar cada figura manualmente antes de compilar el libro.
\usepackage[mode=image, subpreambles=false]{standalone}

% --- 2. MATEMÁTICAS Y CIENCIAS ---
\usepackage{amsmath, amssymb, amsfonts, mathtools}
\usepackage{enumitem}

% LaTeX trae \DeclareMathSizes{5}{5}{5}{5}, así que a 5 pt (\tiny) los subíndices
% salen del mismo tamaño que la base: en las etiquetas de figuras el M de
% $\omega_M$ queda desproporcionado. Se restablece la proporción habitual.
\DeclareMathSizes{5}{5}{3.7}{3}

% --- 3. DISEÑO Y GEOMETRÍA --- 
% Unificamos los márgenes para todo el documento
\usepackage[left=2.5cm,right=2.5cm,top=3cm,bottom=3cm]{geometry}
\makeatletter
\ifstandalone % Evita que geometry dañe el recorte de las figuras bajándolas 3cm
  \@ifclasswith{standalone}{crop=false}{
    % Es un capítulo/sección: Dejamos que geometry actúe.
  }{
    % Es una figura: Neutralizamos geometry para que no las recorte mal.
    \geometry{pass}
  }
\fi
\makeatother
\usepackage{parskip} % Espaciado entre párrafos sin sangría
\usepackage{float}   % Para usar [H] en figura

% Ajustes de flotantes para evitar espacios verticales exagerados
\renewcommand{\topfraction}{0.95}      % Permite que las figuras ocupen hasta el 95% superior
\renewcommand{\bottomfraction}{0.95}   % Permite que las figuras ocupen hasta el 95% inferior
\renewcommand{\textfraction}{0.05}     % Permite hojas con tan solo un 5% de texto
\renewcommand{\floatpagefraction}{0.8} % Requiere que una página de solo figuras esté 80% llena
\setcounter{topnumber}{4}
\setcounter{bottomnumber}{4}
\setcounter{totalnumber}{10}

% --- 4. GRÁFICOS Y TABLAS ---
\usepackage{graphicx}
\usepackage{booktabs} % Tablas profesionales
\usepackage{caption}  % Control de captions y \captionof
\usepackage{tikz}
\usetikzlibrary{arrows.meta, calc, angles, quotes, babel, backgrounds}
\usepackage{pgfplots}
\usepgfplotslibrary{groupplots}
\usepgfplotslibrary{external}
\usepgfplotslibrary{patchplots}
\pgfplotsset{compat=1.18}
\usepackage[american, straightvoltages]{circuitikz}

% --- 5. COLORES Y CAJAS ---

\usepackage[table]{xcolor}
\usepackage{tcolorbox}

% Definición de colores corporativos/estilo
\definecolor{utnblue}{RGB}{0, 51, 102}
\definecolor{codegreen}{rgb}{0,0.6,0}
\definecolor{codegray}{rgb}{0.5,0.5,0.5}
\definecolor{codepurple}{rgb}{0.58,0,0.82}
\definecolor{backcolour}{rgb}{0.96,0.96,0.96}

% --- 6. ENTORNOS PERSONALIZADOS (CAJAS) ---

% Caja "Resultado" (Azul Corporativo) — Definiciones, fórmulas clave, resultados importantes.
% Uso: \begin{resultado}[Fórmula de Euler] ... \end{resultado}
\newtcolorbox{resultado}[1][]{
    colback=utnblue!5!white,
    colframe=utnblue,
    title=\textbf{#1},
    fonttitle=\bfseries,
    sharp corners,
    boxrule=0.5mm
}

% Caja "Nota" (Gris) — Advertencias, aplicaciones, contexto secundario.
% Uso: \begin{nota}[Cuidado con la calculadora] ... \end{nota}
% Uso: \begin{nota} ... \end{nota}  → título por defecto "Nota"
\newtcolorbox{nota}[1][Nota]{
    colback=codegray!5!white,
    colframe=codegray,
    title=\textbf{#1},
    fonttitle=\bfseries,
    sharp corners,
    boxrule=0.5mm
}

% Caja inline para resaltar un valor y enlazarlo con su otra aparición en una tabla
% o en una cuenta: mismo color, mismo valor.
% Uso: \marcavalor{$4$}  o  \marcavalor[codegreen]{$4$}
\newtcbox{\marcavalor}[1][utnblue]{
    on line,
    boxsep=1pt, left=2pt, right=2pt, top=1pt, bottom=1pt,
    colback=#1!10!white,
    colframe=#1,
    boxrule=0.4pt,
    arc=2pt
}

% --- 7. CÓDIGO FUENTE (LISTINGS) ---
\usepackage{listings}

\lstdefinestyle{miestilo}{
    backgroundcolor=\color{backcolour},   
    commentstyle=\color{codegreen},
    keywordstyle=\color{magenta},
    numberstyle=\tiny\color{codegray},
    stringstyle=\color{codepurple},
    basicstyle=\ttfamily\footnotesize,
    breakatwhitespace=false,         
    breaklines=true,                 
    captionpos=b,                    
    keepspaces=true,                 
    numbers=left,                    
    numbersep=5pt,                  
    showspaces=false,                
    showstringspaces=false,
    showtabs=false,                  
    tabsize=2,
    frame=single,                    
    rulecolor=\color{codegray}
}

\lstset{style=miestilo}
% Soporte para tildes y ñ en bloques de código
\lstset{literate={ñ}{{\~n}}1 {á}{{\'a}}1 {é}{{\'e}}1 {í}{{\'i}}1 {ó}{{\'o}}1 {ú}{{\'u}}1}

% Operadores matemáticos personalizados

% Vector en sentido discreto (lista finita de coordenadas): negrita + flecha.
% Uso: \vect{v}, \vect{e}_k, \vect{0}.
\newcommand{\vect}[1]{\vec{\mathbf{#1}}}

% Fecha de versión y comando para capítulos standalone
\providecommand{\verdate}{\the\year.\ifnum\month<10 0\fi\the\month.\ifnum\day<10 0\fi\the\day}
\newcommand{\chapterversion}{%
    \vspace{-1.5cm}\begin{flushright}\small\textit{\color{codegray}Versión~\verdate}\end{flushright}\vspace{0.5cm}%
}

% --- 8. HIPERVÍNCULOS (DEBE SER EL ÚLTIMO PAQUETE) ---
\usepackage[colorlinks=true, linkcolor=utnblue, urlcolor=utnblue, citecolor=utnblue]{hyperref}

% --- 9. REFERENCIAS CRUZADAS ENTRE CAPÍTULOS STANDALONE ---
% Cuando se compila un capítulo standalone, xr-hyper lee los .aux de los
% otros capítulos (si existen) para resolver \ref{} a labels externos.
% En la compilación del libro completo este bloque no se activa.
\ifstandalone
  \usepackage{xr-hyper}
  \makeatletter
  \newcommand{\xrchap}[1]{%
    \edef\xrc@self{\detokenize{Capitulo#1}}%
    \edef\xrc@job{\jobname}%
    \ifx\xrc@self\xrc@job\else
      \IfFileExists{../#1capitulo/Capitulo#1.aux}%
        {\externaldocument{../#1capitulo/Capitulo#1}}{}%
    \fi
  }
  \makeatother
  \xrchap{01}\xrchap{02}\xrchap{03}\xrchap{04}
  \xrchap{05}\xrchap{06}\xrchap{07}\xrchap{08}
  \xrchap{09}\xrchap{10}\xrchap{11}\xrchap{12}
  \xrchap{13}
\fi
% \PassOptionsToPackage{gray}{xcolor}

% fig_replicacion
% Tres paneles apilados sobre el mismo eje de frecuencias:
%   (a) X(jw): triangulo en [-wM, wM], wM=1.
%   (b) espectro del peine: impulsos de peso 2pi/Delta t en cada k*ws, ws=3.
%   (c) X_s(jw) = (a) convolucionada con (b): copias del triangulo en 0, +-ws,
%       +-2ws, todas a escala 1/Delta t (ws=3 > 2 wM=2, no se solapan).

\begin{document}
\begin{tikzpicture}[>={Latex[length=4pt]}, font=\scriptsize, xscale=0.62,
    tri/.style={utnblue, thick, fill=utnblue!12},
    tricopy/.style={utnblue!70, thick, fill=utnblue!7},
    imp/.style={utnblue, thick, -{Latex[length=4pt]}},
    ax/.style={->, thin, black},
]
    % ===================== (a) X(jw) =====================
    \begin{scope}[shift={(0,0)}]
        \node at (-8.6,0.55) {(a)};
        \draw[ax] (-8,0) -- (8,0) node[right] {$\omega$};
        \draw[tri] (-1,0) -- (0,1.0) -- (1,0) -- cycle;
        \node[utnblue, anchor=south] at (0,1.02) {$X(j\omega)$};
        % ticks +-wM
        \draw[thin] (1,0.06) -- (1,-0.06) node[below, font=\tiny] {$\omega_M$};
        \draw[thin] (-1,0.06) -- (-1,-0.06) node[below, font=\tiny] {$-\omega_M$};
    \end{scope}

    % ===================== (b) espectro del peine =====================
    \begin{scope}[shift={(0,-2.7)}]
        \node at (-8.6,0.55) {(b)};
        \draw[ax] (-8,0) -- (8,0) node[right] {$\omega$};
        % impulsos en 0, +-ws, +-2ws
        \draw[imp] (-6,0) -- (-6,0.85);
        \draw[imp] (-3,0) -- (-3,0.85);
        \draw[imp] ( 0,0) -- ( 0,0.85);
        \draw[imp] ( 3,0) -- ( 3,0.85);
        \draw[imp] ( 6,0) -- ( 6,0.85);
        % puntos suspensivos: los impulsos siguen sin fin
        \node[utnblue, font=\tiny] at (-7.4,0.42) {$\cdots$};
        \node[utnblue, font=\tiny] at ( 7.4,0.42) {$\cdots$};
        % peso de cada impulso, anotado sobre el del origen
        \node[anchor=west, font=\tiny] at (0.18,0.86) {$\tfrac{2\pi}{\Delta t}$};
        % ticks en multiplos de ws
        \draw[thin] (3,0.06) -- (3,-0.06) node[below, font=\tiny] {$\omega_s$};
        \draw[thin] (6,0.06) -- (6,-0.06) node[below, font=\tiny] {$2\omega_s$};
        \draw[thin] (-3,0.06) -- (-3,-0.06) node[below, font=\tiny] {$-\omega_s$};
        \draw[thin] (-6,0.06) -- (-6,-0.06) node[below, font=\tiny] {$-2\omega_s$};
    \end{scope}

    % ===================== (c) X_s(jw) =====================
    \begin{scope}[shift={(0,-5.4)}]
        \node at (-8.6,0.55) {(c)};
        \draw[ax] (-8,0) -- (8,0) node[right] {$\omega$};
        % copias laterales
        \draw[tricopy] (-7,0) -- (-6,0.85) -- (-5,0) -- cycle;
        \draw[tricopy] (-4,0) -- (-3,0.85) -- (-2,0) -- cycle;
        \draw[tricopy] ( 2,0) -- ( 3,0.85) -- ( 4,0) -- cycle;
        \draw[tricopy] ( 5,0) -- ( 6,0.85) -- ( 7,0) -- cycle;
        % copia central (destacada)
        \draw[tri] (-1,0) -- (0,0.85) -- (1,0) -- cycle;
        \node[utnblue, anchor=south] at (0,1.20) {$X_s(j\omega)$};
        % puntos suspensivos: las copias siguen sin fin
        \node[utnblue!70, font=\tiny] at (-7.6,0.42) {$\cdots$};
        \node[utnblue!70, font=\tiny] at ( 7.6,0.42) {$\cdots$};
        % eje y: marca de la escala 1/Delta t
        \draw[ax] (0,-0.05) -- (0,1.15);
        \draw[thin] (-0.12,0.85) -- (0.12,0.85);
        \node[anchor=east, font=\tiny] at (-0.18,0.92) {$\tfrac{1}{\Delta t}$};
        % ticks en multiplos de ws
        \draw[thin] (3,0.06) -- (3,-0.06) node[below, font=\tiny] {$\omega_s$};
        \draw[thin] (6,0.06) -- (6,-0.06) node[below, font=\tiny] {$2\omega_s$};
        \draw[thin] (-3,0.06) -- (-3,-0.06) node[below, font=\tiny] {$-\omega_s$};
        \draw[thin] (-6,0.06) -- (-6,-0.06) node[below, font=\tiny] {$-2\omega_s$};
        % marca de la separación entre copias
        \draw[<->, gray!65] (1,-0.62) -- (2,-0.62);
        \node[gray!65, anchor=north, font=\tiny] at (1.5,-0.60) {separación};
    \end{scope}
\end{tikzpicture}
\end{document}
